\chapter{Desenvolvimento}

Este capítulo apresenta os sistemas desenvolvidos pela Faculdade para dar suporte aos processos administrativos e acadêmicos, implementados com o uso do framework \textit{Laravel}. Também são contempladas as bibliotecas criadas em PHP com o objetivo de tornar o processo de reutilização de código mais eficiente entre os sistemas.

É fundamental que os responsáveis tenham atenção especial às bibliotecas e aos sistemas que estão no âmbito da organização \texttt{uspdev}, pois esses projetos exigem articulação com os demais desenvolvedores da Universidade. Já os projetos sob o escopo \texttt{fflch} permitem maior autonomia para modificações e evoluções.

Para diferenciar os projetos, deve-se observar o \textit{namespace} do repositório. Por exemplo: \texttt{uspdev/projeto1} indica que o projeto pertence ao \texttt{uspdev}, enquanto \texttt{fflch/projeto2} pertence à FFLCH.

Em ambos os casos, é obrigatório realizar um \textit{fork} do projeto, evitando trabalhar diretamente no repositório \textit{upstream}.

Os responsáveis também devem analisar cuidadosamente o código enviado pela comunidade do \texttt{uspdev}, interagindo com a comunidade, quando necessário, assim como o código enviado pelos nossos estagiários por meio de \textit{pull requests}, garantindo qualidade, padronização e aderência às boas práticas de desenvolvimento.

\section{Sistemas}
\loadjson{dev}

\newpage
\section{Bibliotecas}
\loadjson{bib}

%\textbf{Sistemas previstos para entrada em produção:}
%\begin{itemize}
%    \item \texttt{uspdev/patrimonio-complementar}
%    \item \texttt{uspdev/datagrad}
%    \item \texttt{uspdev/solicitacoes-documentos}
%    \item \texttt{uspdev/evasao}
%    \item \texttt{uspdev/portal-sistemas}
%\end{itemize}